\section{\blue{Related Work}}
\label{sec:related}

\blue{
The shape, the level and the amplitude of a periodic series have each been modeled before, but in separate lines of work. Decomposition separates a trend, tokens and cross-cycle models capture the cycle shape, and normalization removes and restores the level and the scale; in none of them does the amplitude have a token of its own, driven by the level. We review each line by what it covers. Where a line leaves two kinds of information mixed, its network must grow to separate them; where it leaves one out, its forecast loses accuracy. \sysplus covers all three with its tokens, and this coverage is what lets it be small. We close with forecasting at the scale of a sensor network, where this matters most.

\fakeparagraph{Decomposition}
Decomposing a series into interpretable parts is a long tradition. Classical methods separate a trend, a seasonal component and a remainder, by local regression in STL~\citep{cleveland1990stl} or by exponential smoothing in Holt--Winters forecasting~\citep{holt2004forecasting,winters1960forecasting,hyndman2021fpp3}, whose multiplicative form scales the seasonal component with the level. Deep forecasters build the split into their architecture. Autoformer~\citep{wu2021autoformer} and FEDformer~\citep{zhou2022fedformer} separate a moving-average trend, a mixture of moving averages in FEDformer, from a seasonal part inside every layer, and model the seasonal part with auto-correlation or frequency-domain attention. DLinear~\citep{zeng2023dlinear} applies one linear map to each part. TimeMixer~\citep{timemixer24} downsamples the series to several scales, decomposes each scale into seasonal and trend parts, mixes the seasonal parts from fine to coarse and the trend parts from coarse to fine, and sums the forecasts of one predictor per scale. The interpretable N-BEATS~\citep{oreshkin2020nbeats} expands the forecast in polynomial trend and Fourier seasonal bases, and CycleNet~\citep{lin2024cyclenet} learns one recurrent cycle per channel, fixed after training, and forecasts the residual with a separate model. These designs share the premise of this paper, that the parts of a series should not be learned in one undifferentiated representation.

\fakeparagraph{Where decomposition falls short}
In the terms of the form $X\approx\ell\,\one^\top+a\,s^\top$ of Section~\ref{sec:mot_tokens}, the trend of these deep decompositions stands in for the level, the seasonal part or the learned cycle holds the shape, and the amplitude of each cycle is left inside the seasonal part or the residual, with no component of its own. They therefore cover the three kinds of information only partly, in three respects. First, they split the series along time, not into its cycles: a level that one number per cycle describes is carried by one value per time step, which the trend model must forecast step by step. Even the coarsest scale of TimeMixer comes from downsampling by fixed factors that the period does not set, so its trend remains a curve over time rather than one level per cycle. Second, they add the parts, and nothing ties the seasonal swing to the trend, so the seasonal model must learn the coupling from examples, level by level; TimeMixer adds its parts twice, when it recombines the seasonal and trend parts at each scale and when it sums the forecasts of its per-scale predictors. The multiplicative Holt--Winters model ties the swing to the level, but in a fixed proportion and through a per-series recursion with a few smoothing constants, not through a coupling learned by a forecaster that serves every sensor. Third, the network that receives each part must untangle whatever the split leaves mixed, and it grows accordingly; where it is a single linear map, as in DLinear, it cannot separate what stays mixed, and the forecast pays in accuracy instead. Partial coverage thus forces a choice between capacity and accuracy.

\sysplus instead represents the folded window by one token for each kind of information: one level per cycle, an amplitude driven by the level, and phase tokens that carry the shape. Its level pathway combines a few fixed horizon profiles, as the trend basis of N-BEATS combines low-degree polynomials, and its amplitude pathway scales the cycle shape with the level while provably leaving every future cycle mean unchanged (Proposition~\ref{prop:modulation}). Each pathway forecasts one kind of information in its low-dimensional form and stays small (Section~\ref{sec:theory}). Where TimeMixer adds scales and predictors, \sysplus adds one number per cycle and a level-driven amplitude factor: it is the full coverage of the shape, the level and the amplitude, not a restriction of the model, that lets it be small.

\fakeparagraph{Tokens and cross-cycle models}
The token decides what a forecaster must untangle. Autoformer and FEDformer embed every time step as a token. PatchTST~\citep{nie2023patchtst}, Crossformer~\citep{zhang2023crossformer} and PatchMLP~\citep{tang2025patchmlp} take patches of consecutive values, which mix the shape, the level and the amplitude of a cycle, so that deep attention or MLP stacks must separate them; Crossformer also exchanges information across variables through a few router tokens, and the routers of \sysplus, a few learned tokens that gather messages from all phases and distribute them back at a cost linear in the number of phases, follow these and the latent array through which Perceiver~\citep{jaegle2021perceiver} reads a large input. iTransformer~\citep{liu2023itransformer} takes the whole series of a variable. TimesNet~\citep{wu2023timesnet} folds a series by each dominant period into a two-dimensional array, as \sysplus does, but processes the arrays of several periods with stacked two-dimensional convolution blocks, which must themselves separate the level and the amplitude from the shape; once the level and the amplitude are represented, a single period already carries 78\% to 93\% of the energy of a normalized window (Fig.~\ref{fig:token_energy}).

Cross-cycle models use the fold to stay small. Among them, SparseTSF~\citep{lin2024sparsetsf} forecasts each phase with one shared linear map, and TimeBase~\citep{huang2025timebase} forecasts a few basis components of cycle-length segments. Neither gives the level or the amplitude a token of its own. Both are linear in the normalized window and read the level and the shape with the same weights over the cycles: a long average gives a clean shape but a stale level, and a short one a current level but a noisy shape. Neither can therefore average the shape over the window and scale it by a level forecast from the recent cycles, a product of two different reads of the window, which the amplitude pathway of \sysplus forms (Proposition~\ref{prop:modulation}). Beyond the low rank that TimeBase exploits, none of these works separates what a token must represent into the shape, the level and the amplitude of the cycles. \sysplus does, and its analysis proves when the tokens obtained by grouping values across cycles stay in a fixed subspace (Theorem~\ref{thm:stability}) and bounds what a forecaster that reads them through a layer normalization can produce (Section~\ref{sec:theory}).

\fakeparagraph{Normalization: the level and the scale as statistics}
RevIN~\citep{kim2021reversible}, which \sysplus uses, subtracts the mean of each input window, divides it by its standard deviation, and inverts both at the output. The Non-stationary Transformer~\citep{nonstationarytrans2022} reinjects factors learned from the removed statistics into attention, and ST-Norm~\citep{deng2021st} normalizes along the temporal and spatial dimensions of a sensor network. Dish-TS~\citep{fan2023dishts}, SAN~\citep{liu2023san} and FAN~\citep{ye2024fan} forecast the statistics of the output window: its level and scale, the mean and standard deviation of each slice, and its dominant frequency components. These plug-in methods treat the non-stationary statistics as quantities to remove before a forecaster and to restore after it; SAN, whose slices can span one cycle, comes closest to the level and the amplitude, but forecasts both with a statistics module outside the forecaster. \sysplus treats the level and the amplitude as information inside the forecaster. Normalizing a window removes one mean and one standard deviation, but the level that drifts across the cycles of the normalized window and the amplitude that moves with it remain, and the level alone carries up to 80\% of the energy of a normalized window (Fig.~\ref{fig:token_energy}). The level and the amplitude of the cycles therefore have tokens of their own, and the amplitude is driven by the level forecast rather than forecast on its own. The envelopes of Section~\ref{sec:theory} show why this matters: a forecaster whose affine predictor reads the output of a layer normalization, as the phase pathway does, confines the forecast level and amplitude, for any fixed weights, to envelopes that no added depth, routers or heads remove. Because the envelopes are stated in normalized coordinates, de-normalization rescales them but does not lift them; normalization and level tokens do different work.

\fakeparagraph{Forecasting at the scale of a sensor network}
Spatio-temporal graph models such as Graph WaveNet~\citep{wu2019graphwavenet}, MTGNN~\citep{wu2020mtgnn}, AGCRN~\citep{bai2020agcrn} and D2STGNN~\citep{shao2022decoupled} pass messages along a given or learned graph over the sensors, and AutoCTS~\citep{wu2021autocts} searches such architectures automatically. Other work lowers their cost: BigST~\citep{han2024bigst} makes message passing linear in the number of sensors, LightCTS~\citep{lai2023lightcts} uses lightweight temporal and spatial operators, and STID~\citep{shao2022stid} replaces the graph with spatial and temporal identity embeddings and an MLP. \sysplus is complementary to this line: it shrinks the channel-independent model that serves every sensor, and a spatial model can be built on top of it. A second route to efficiency keeps the model applied to each series small: DLinear~\citep{zeng2023dlinear} and FITS~\citep{fits2024} with linear maps, FilterNet~\citep{filternet2024} with learnable frequency-domain filters, and SSCNN~\citep{deng2024parsimony} with a parsimonious decomposition. A linear map keeps a model small but cannot form the level-driven amplitude, as discussed above for the cross-cycle models. \sysplus is small by coverage instead: its tokens give each pathway one low-dimensional part of the window to forecast, and each pathway is restricted only to the form that its part takes. The structure it exploits is, moreover, at its strongest on sensor networks: the amplitude follows the level most closely on the 862 road sensors of Traffic and the 321 meters of Electricity (Section~\ref{sec:mot_changes}), where \sysplus also leads SparseTSF and TimeBase, which fold the series as it does, by about 5\% to 6\% (Section~\ref{sec:exp_results}). Complete tokens pay off most where one model serves the most sensors.
}
